% II. Анализ на възможните решения и обосновка на избора.
% Три независими избора: метод за многокритериална оценка, метод за решаване на
% задачата за подбор, технологична основа. Всеки се решава отделно.
\section{Анализ на възможните решения и обосновка на избора}
\label{sec:alternatives}

\subsection{Метод за многокритериална оценка}

Разгледани са два метода. Методът на аналитичната йерархия~\cite{saaty1980ahp}
извежда теглата на критериите от попарни сравнения и дава на оценяващия мярка за
съгласуваност на собствените му преценки. Цената е обемът на въпросите: за $n$
критерия сравненията са $n(n-1)/2$, а при добавяне на нов критерий целият набор
се преразглежда. Методът TOPSIS~\cite{hwang1981topsis} изисква теглата да бъдат
дадени наготово и класира алтернативите по относителна близост до идеалната
точка, което го прави евтин при много алтернативи и стабилен при добавянето на
нова.

Изборът е двата метода да бъдат реализирани, а не единият. Причината е емпирична
и е посочена в прегледа: върху едни и същи данни различните методи могат да
класират различно~\cite{triantaphyllou2000mcdm}. Реализирането само на единия би
скрило този факт зад числото, което той произвежда. При двата реализирани метода
несъвпадението в класирането става наблюдаемо и се превръща в резултат на
проекта, а не в мълчаливо допускане. Аналитичната йерархия се използва за
извличане на теглата, а TOPSIS за класирането при тези тегла, с възможност и
двете да се сменят.

\subsection{Метод за решаване на задачата за подбор}

Задачата за избор на портфейл при бюджет, капацитет по периоди и отношения на
предхождане е формулирана като задача на целочисленото линейно програмиране.
Разгледани са три пътя за решаването ѝ.

\begin{enumerate}[topsep=0pt,itemsep=0pt,leftmargin=*]
\item \textbf{Лакома евристика} по съотношение полза към стойност. Евтина за
      реализация, но не дава гаранция за оптималност и не се справя с
      ограниченията за предхождане без допълнителни поправки.
\item \textbf{Собствен алгоритъм разклоняване и ограничаване} по схемата
      от~\cite{martello1990knapsack}. Дава точно решение и е напълно прозрачен,
      но за многомерния вариант с предхождане изисква съществено повече работа,
      отколкото проектът предвижда за този слой.
\item \textbf{Готов решател}, конкретно решателят за задачи с ограничения от
      пакета OR-Tools. Дава точно решение, поема многомерните ограничения без
      допълнителна работа и позволява същият модел да се преизползва при анализа
      на чувствителността.
\end{enumerate}

Избран е третият път. Основанието е разпределението на усилието: приносът на
проекта е в моделирането и в проследимостта на препоръката, не в реализацията на
решател. Лакомата евристика остава реализирана като контролна стойност, срещу
която се проверява дали решателят връща поне толкова добро решение. Това е и
единственият начин грешка в модела да бъде забелязана: решател, който връща
по-лош резултат от евристиката, не е бавен, а сгрешен.

\subsection{Технологична основа}

Ядрото на модела е реализирано на C++20. Изборът е обоснован от два фактора:
анализът на чувствителността изпълнява един и същ модел стотици пъти при
изменени входове, което поставя изисквания към времето за изпълнение, а пакетът
OR-Tools предоставя първокласен интерфейс на същия език. Алтернативата, Python с
интерфейса на същия пакет, би съкратила времето за писане, но би довела до втори
език в системата без друга полза.

Потребителският интерфейс е отделен от ядрото и общува с него чрез описание на
сценария в JSON. Границата е поставена там съзнателно: сценарият е файл, който
може да се версионира, да се сравнява и да се прилага към проекта като
доказателство какви са били допусканията. Интерфейсът е уеб базиран, за да не
въвежда зависимост от платформен инструментариум за изобразяване.
\TODO{потвърди окончателния избор на HTTP слой и го запиши тук; към момента на
писане решението не е фиксирано}
